% =====================================================================
%  Section 1.1 "Project specifications and requirements" for the master
%  report (Chapter 1, "The Problem: Visualising the Invisible").
%  Drop the body (from \section onward) into the master document; the
%  preamble here is only so it previews standalone.
%  Requirement statements are quoted from the EE2 "maths accelerator" brief:
%  https://github.com/edstott/EE2Project/blob/main/maths-accelerator/README.md
% =====================================================================
\documentclass[11pt]{report}
\usepackage[margin=1in]{geometry}
\usepackage{array}
\usepackage{booktabs}
\begin{document}

\section{Project specifications and requirements}

The project brief asks for a \emph{real-time visualisation of a computationally
intensive, embarrassingly-parallel mathematical function}, computed on the
supplied PYNQ--Z1 board with a hardware accelerator for the inner loops, compared
against a CPU implementation, and wrapped in a user interface that makes it
useful as an educational tool. The brief suggests fractals as a typical choice
and explicitly discourages N-body simulations, because their global state and
random memory-access patterns map poorly onto the FPGA.

We chose a more ambitious ``function'': the time-dependent solution of Maxwell's
equations by the Finite-Difference Time-Domain (FDTD) method, visualised either
as a 2D colour field or as a ray-marched 3D surface. This makes the otherwise
invisible electromagnetic field tangible, which is a stronger educational result
than a synthetic fractal, while still satisfying the brief's computational
demands. The choice does change the shape of the parallelism, and it is worth
being precise about how:

\begin{itemize}
  \item \textbf{The renderer is embarrassingly parallel.} Each of the
  $640\times480$ output pixels is produced by an independent ray-march with no
  data dependence on any other pixel --- exactly the property the brief asks
  for, and the source of most of the throughput demand. This is the part that
  must keep up with the display.
  \item \textbf{The FDTD field has only local, structured dependence.} Updating a
  cell needs only its immediate neighbours, and every memory access is regular
  and known at compile time. This is the favourable opposite of the discouraged
  N-body case (global coupling, random access): the locality is precisely what
  lets us stream and spatially parallelise the solver on the fabric (the
  four-lane decomposition in the FDTD chapter).
\end{itemize}

Table~\ref{tab:reqs} lists the brief's requirements and where each is addressed
in this report; the evidence that they were met (throughput numbers, timing,
resource use) is collected in the Evaluation chapter.

\begin{table}[h]
\centering
\renewcommand{\arraystretch}{1.25}
\begin{tabular}{|p{0.46\textwidth}|p{0.46\textwidth}|}
\hline
\textbf{Requirement (brief)} & \textbf{How the project addresses it} \\
\hline
\multicolumn{2}{|l|}{\textit{Visualisation}}\\
\hline
Display a visualisation of a mathematical function computed in real time. &
The FDTD solution of Maxwell's equations is rendered to HDMI at $640\times480$,
25\,fps, updating live (Ch.~3--5). \\
\hline
The function shall be computationally intensive --- non-trivial at the required
resolution and frame rate. &
A $128\times128$ field is time-stepped every frame and ray-marched per pixel; the
per-pixel renderer alone is the binding throughput constraint (Ch.~4). \\
\hline
The computation should be embarrassingly parallel (pixels independent). &
The renderer ray-marches every pixel independently with no inter-pixel data
dependence; the FDTD field has only local nearest-neighbour coupling and is
parallelised by spatial decomposition (Ch.~3, Ch.~4). \\
\hline
\multicolumn{2}{|l|}{\textit{Accelerator}}\\
\hline
Generated on the PYNQ--Z1, with an accelerator in FPGA soft logic computing the
inner loops. &
The FDTD update and the ray-march inner loops both run in the programmable logic;
the processor only configures and reads back (Ch.~3--5). \\
\hline
The accelerator shall be described in Verilog or SystemVerilog. &
All datapath RTL (solver, magnitude, bridge, renderer) is SystemVerilog/Verilog,
built from a scripted, version-controlled flow (Ch.~5). \\
\hline
The accelerator shall provide a CPU interface for parameter adjustment. &
An AXI-GPIO register map exposes amplitude, frequency, source position, display
mode, motion and camera to the processor, driven live from a notebook (Ch.~5). \\
\hline
Number formats and word lengths chosen to trade accuracy against throughput. &
Signed Q3.13 fixed-point throughout, selected from a documented comparison of
formats (Q1.15--FP32) against rounding error and resource cost (Ch.~3). \\
\hline
Throughput shall exceed a CPU-only C/C++/Cython implementation. &
Benchmarked against an ARM (NEON) software baseline; the accelerated solver is
substantially faster (Ch.~3, Ch.~7). \\
\hline
\multicolumn{2}{|l|}{\textit{User interface}}\\
\hline
Provide a user interface that enhances educational value. &
The user can place and move sources, switch field/energy views, vary speed and
material, and watch propagation, interference, the Doppler effect and energy flow
in real time (Ch.~6). \\
\hline
The UI may use separate hardware. &
A standalone ESP32 with a resistive touch panel and an analog conducting-sheet
rig form the input hardware, linked to the board over serial (Ch.~6). \\
\hline
The UI shall allow intuitive or interesting parameter adjustment. &
The source is placed by dragging a stylus on the panel; potentiometers set
amplitude, material and camera; the conducting-sheet probe links the simulation
to a measurable physical analogue (Ch.~6). \\
\hline
The UI should present information as an overlay or a separate medium. &
Parameters and status are shown on the host notebook (separate medium), the
physical panel is labelled, and a colour-vision-safe palette communicates field
magnitude (Ch.~6). \\
\hline
\end{tabular}
\caption{Project requirements and where each is addressed. Quantitative evidence
that they were met is given in the Evaluation chapter.}
\label{tab:reqs}
\end{table}

\end{document}
